\section{Our Go Implementation}


The entire lattice cryptography stack for Lux Network is implemented in Go with SIMD
acceleration via assembly (covered in our companion NTT paper).

\subsection{Package Structure}

\begin{lstlisting}[language=bash,caption={Package layout}]
luxfi/crypto/
    lattice/
        ring.go       # Ring element operations
        ntt.go        # NTT forward/inverse
        ntt_amd64.s   # AVX2 NTT butterfly
        poly.go       # Polynomial arithmetic
        sample.go     # Sampling (CBD, uniform, challenge)
    mlkem/
        keygen.go     # ML-KEM key generation
        encaps.go     # Encapsulation
        decaps.go     # Decapsulation
    mldsa/
        keygen.go     # ML-DSA key generation
        sign.go       # Signing with rejection sampling
        verify.go     # Verification
    ringtail/
        threshold.go  # Threshold operations
        share.go      # Secret share management
        aggregate.go  # Partial signature aggregation
\end{lstlisting}

\subsection{Performance}

\begin{table}[H]
\centering
\small
\begin{tabular}{@{}lrrl@{}}
\toprule
\textbf{Operation} & \textbf{Time} & \textbf{vs.\ Reference} & \textbf{Platform} \\
\midrule
ML-KEM-768 KeyGen & 28\,$\mu$s & 1.1x ref C & Apple M2 \\
ML-KEM-768 Encaps & 34\,$\mu$s & 1.2x ref C & Apple M2 \\
ML-KEM-768 Decaps & 38\,$\mu$s & 1.1x ref C & Apple M2 \\
ML-DSA-65 KeyGen & 52\,$\mu$s & 1.3x ref C & Apple M2 \\
ML-DSA-65 Sign & 186\,$\mu$s & 1.4x ref C & Apple M2 \\
ML-DSA-65 Verify & 56\,$\mu$s & 1.2x ref C & Apple M2 \\
NTT-256 (pure Go) & 3.1\,$\mu$s & 2.8x ref C & Apple M2 \\
NTT-256 (AVX2) & 0.9\,$\mu$s & 0.95x ref C & AMD Zen 4 \\
\bottomrule
\end{tabular}
\caption{Performance of Lux Network's Go lattice crypto implementation.}
\end{table}

%% ============================================================
